\subsection{Implementation plan}

The implementation of BBP will be executed in a phased manner to manage the complexities of a distributed system. 

Initially, development will focus on the infrastructure layer using Spring Cloud and Netflix components. 
This includes setting up the Service Discovery and API Gateway to establish a robust routing mechanism for all incoming client requests.

Once the infrastructure is stable, work will proceed to the User Service to manage registration and authentication processes. 

Following this, the Trip Service and Path Service will be developed to handle the core logic of recording cycling activities and computing path scores based on user-submitted status reports. 

Finally, the Data Processing Service will be implemented to consume events from the Kafka broker, allowing for the generation of performance statistics and the integration of external meteorological data.

\subsection{Integration plan}

The integration strategy for BBP combines synchronous and asynchronous patterns to balance performance with reliability. Direct interactions, such as user logins or path searches, will be handled synchronously through the API Gateway via RESTful calls, given the Microservices architecture. This ensures immediate feedback for high-priority user actions.

Instead more complex operations will utilize an event-driven approach. When the Trip Service completes a recording, it will publish a "Trip Recorded" event to the Kafka message broker. This allows the Data Processing Service to aggregate statistics and the Path Service to update the path database without blocking the user's mobile application. 

Furthermore, external integrations with Map and Weather providers will be encapsulated within their respective services to maintain the principle of loose coupling of Microservices. 

\subsection{Test plan}

The testing phase ensures that all functional requirements are met and that the system remains stable under load in each development stage and for each microservice within the architecture.

\begin{itemize}
   \item \textbf{Unit Testing:} Each one of the microservices will undergo rigorous unit testing using JUnit to validate business logic in isolation.
   \item \textbf{Integration Testing:} We will verify the exchange of messages through a tool like Kafka and ensure that the three independent databases (User, Trip, and Path) maintain data integrity during each type of operations.
   \item \textbf{Functional and System Testing} This will focus on critical BBP scenarios, such as the accuracy of path score calculations and the reliability of the automated mode in reconstructing GPS paths. \newline We will also simulate concurrent user access to ensure the system can handle the expected community load.
   \item \textbf{User Interface Testing:} The mobile application will be tested on various devices and operating systems to ensure a consistent user experience and responsiveness and ease of use from client perspective.
   \item \textbf{Security Testing:} We will conduct vulnerability assessments and penetration testing to safeguard user data, particularly focusing on authentication and authorization mechanisms, utilizing notorious hacking strategies.
   \item \textbf{Performance Testing:} Load testing will be conducted to assess system performance under peak conditions, ensuring that response time remain acceptable.
\end{itemize}

